\documentclass[../main.tex]{subfiles}
\begin{document}

Trong bài toán điều tra số và phát hiện tấn công trên hệ điều hành \textbf{Windows}, việc chỉ dựa vào các luật phát hiện tĩnh hoặc các chỉ báo đơn lẻ thường không đủ để mô tả đầy đủ mức độ nguy hiểm của một tiến trình. Các hành vi tấn công hiện đại, đặc biệt là các kỹ thuật \textit{Living-off-the-Land} (LotL) và \textit{Advanced Persistent Threat} (APT), thường được che giấu dưới dạng chuỗi hành vi phân tán, có mức độ bất thường thấp nếu quan sát riêng lẻ từng sự kiện.

Do đó, một hệ thống \textit{triage} hiện đại cần có khả năng tổng hợp nhiều nguồn tín hiệu khác nhau, bao gồm hành vi tiến trình, quan hệ ngữ cảnh, cấu trúc chuỗi thực thi, kết nối mạng và các luật phát hiện đã biết, từ đó lượng hóa mức độ rủi ro dưới dạng điểm số có thể giải thích được (\textit{explainable scoring}). Điều này đòi hỏi sự kết hợp giữa các kỹ thuật thống kê \textit{robust}, học máy không giám sát (\textit{unsupervised learning}), lý thuyết đồ thị (\textit{graph theory}) và mô hình xác suất.

Chương này trình bày các cơ sở lý thuyết và các mô hình toán học được sử dụng trong hệ thống đề xuất. Nội dung chương bao gồm nền tảng về \textbf{Windows Event Logging} và \textbf{Sysmon}, mô hình phát hiện bất thường (\textit{anomaly detection}) theo phân loại của \textbf{Varun Chandola}, các kỹ thuật chuẩn hóa thống kê như \textbf{MAD} và \textbf{NIF}, phương pháp phân cụm \textbf{HDBSCAN}, mô hình \textbf{Markov} cho phân tích chuỗi tiến trình, cũng như cơ chế định lượng và lan truyền rủi ro (\textit{risk propagation}) trên đồ thị tiến trình.

Các công thức và mô hình được trình bày trong chương này đóng vai trò nền tảng cho kiến trúc phân tích và hệ thống thực nghiệm sẽ được mô tả trong các chương tiếp theo.

\section{Nền tảng hệ điều hành Windows và Nhật ký sự kiện (Windows Event Log)}
\label{section:2.1}
Hệ điều hành Windows là một trong những nền tảng phổ biến nhất trong môi trường doanh nghiệp, do đó cũng trở thành mục tiêu chính của các cuộc tấn công mạng. Để phục vụ việc giám sát và điều tra, Windows cung cấp hệ thống ghi log tập trung thông qua Windows Event Log, cho phép lưu trữ và truy xuất các sự kiện liên quan đến hệ thống, bảo mật và ứng dụng.

Windows Event Log được tổ chức thành các kênh log chính như Security, System và Application. Trong đó, log Security đóng vai trò quan trọng trong việc theo dõi các sự kiện liên quan đến xác thực, truy cập và thay đổi quyền hạn. Tuy nhiên, các log mặc định thường thiếu chi tiết về hành vi tiến trình, đặc biệt là các tham số dòng lệnh hoặc quan hệ giữa các tiến trình.

Để khắc phục hạn chế này, công cụ Sysmon (System Monitor) được sử dụng như một nguồn log bổ sung, cung cấp thông tin chi tiết hơn về hoạt động của hệ thống. Theo tài liệu của Microsoft, Sysmon có thể ghi nhận các sự kiện như tạo tiến trình, kết nối mạng, thay đổi registry và truy cập file, giúp nâng cao khả năng phát hiện các hành vi bất thường.

\subsection{Windows Event Logging}
\label{section:2.1.1}
\textbf{Windows} sử dụng cơ chế \textbf{Event Logging} để ghi nhận các hoạt động xảy ra trong hệ thống. Các sự kiện này được lưu dưới dạng file \textit{.evtx}, bao gồm thông tin về tiến trình, mạng, đăng nhập, thay đổi \textit{registry} và nhiều hoạt động khác.

Mỗi sự kiện trong \textbf{Windows Event Log} thường gồm:
\begin{itemize}
    \item \textbf{Timestamp}: Thời gian xảy ra,
    \item \textbf{Event ID} (EID): Mã sự kiện,
    \item \textbf{Provider}: Nguồn sinh log,
    \item \textbf{Level}: Mức độ ưu tiên,
    \item \textbf{Event Data}: Tập dữ liệu mô tả chi tiết hành vi.
\end{itemize}

Các log mặc định của \textbf{Windows} tuy hữu ích cho quản trị hệ thống nhưng còn hạn chế đối với bài toán điều tra số (\textit{digital forensics}). Nhiều hành vi quan trọng như \textit{command line}, \textit{injection} hoặc kết nối mạng của tiến trình thường không được ghi nhận đầy đủ.

Để giải quyết vấn đề này, \textbf{Microsoft} phát triển \textbf{Sysmon} (\textit{System Monitor}) như một thành phần mở rộng phục vụ giám sát hệ thống ở mức chi tiết hơn.

\subsection{Cấu trúc nhật ký sự kiện và vai trò của Sysmon}
\label{subsection:2.1.2}
\textbf{Sysmon} (\textit{System Monitor}) là công cụ thuộc bộ \textbf{Sysinternals} của \textbf{Microsoft}, cho phép ghi nhận chi tiết hoạt động của tiến trình, kết nối mạng và thay đổi hệ thống.

Khác với \textbf{Windows Event Log} thông thường, \textbf{Sysmon} có khả năng:
\begin{itemize}
    \item Ghi \textit{command line} đầy đủ,
    \item Theo dõi quan hệ tiến trình cha–con (\textit{parent-child relationship}),
    \item Ghi nhận \textit{network connection} theo tiến trình,
    \item Phát hiện \textit{injection},
    \item Và lưu nhiều thông tin phục vụ \textit{forensic}.
\end{itemize}

Trong bài toán \textbf{DFIR} \textit{triage}, \textbf{Sysmon} đóng vai trò đặc biệt quan trọng do cung cấp khả năng tái dựng hành vi của tiến trình theo dòng thời gian (\textit{timeline}). Hệ thống trong đề tài sử dụng 9 \textbf{Event ID} chính:

\begin{table}[h]
\centering
\begin{tabular}{|l|l|}
\hline
\textbf{Event ID} & \textbf{Ý nghĩa} \\ \hline
\textbf{1} & \textit{Process Create} \\ \hline
\textbf{3} & \textit{Network Connection} \\ \hline
\textbf{8} & \textit{CreateRemoteThread} \\ \hline
\textbf{10} & \textit{Process Access} \\ \hline
\textbf{11} & \textit{File Create} \\ \hline
\textbf{13} & \textit{Registry Value Set} \\ \hline
\textbf{17} & \textit{Pipe Created} \\ \hline
\textbf{18} & \textit{Pipe Connected} \\ \hline
\textbf{22} & \textit{DNS Query} \\ \hline
\end{tabular}
\caption{Các Event ID chính của Sysmon được sử dụng trong hệ thống}
\end{table}

Trong đó:
\begin{itemize}
    \item \textbf{Event ID 1} là nguồn dữ liệu trung tâm để xây dựng cây tiến trình (\textit{process tree}),
    \item \textbf{Event ID 3} phục vụ phát hiện \textit{beaconing} và \textbf{C2},
    \item \textbf{Event ID 8} và \textbf{10} liên quan đến \textit{injection} và \textit{credential dumping},
    \item \textbf{Event ID 13} phản ánh hành vi \textit{persistence},
    \item \textbf{Event ID 22} hỗ trợ phân tích \textit{DNS-based C2}.
\end{itemize}

\subsection{Kiến trúc tiến trình Windows và vấn đề PID reuse}
\label{subsection:2.1.3}
Trong \textbf{Windows}, mỗi tiến trình được gán một \textbf{Process ID} (PID). Tuy nhiên, \textbf{PID} chỉ là định danh tạm thời và có thể bị tái sử dụng (\textit{PID reuse}) sau khi tiến trình kết thúc.

Ví dụ:
\begin{itemize}
    \item Một tiến trình có \textbf{PID = 4120} kết thúc,
    \item Sau đó \textbf{Windows} có thể cấp lại \textbf{PID = 4120} cho một tiến trình hoàn toàn khác.
\end{itemize}

Điều này tạo ra vấn đề nghiêm trọng trong \textit{forensic} nếu chỉ sử dụng \textbf{PID} để liên kết các sự kiện. Để giải quyết vấn đề này, \textbf{Sysmon} cung cấp \textbf{ProcessGuid}.

\textbf{ProcessGuid} là một định danh duy nhất (\textit{globally unique identifier}) cho mỗi vòng đời tiến trình. Không giống \textbf{PID}, \textbf{ProcessGuid} không bị tái sử dụng và cho phép tái dựng chính xác cây tiến trình (\textit{process tree}).

Do đó, trong hệ thống đề xuất:
\begin{itemize}
    \item \textbf{PID} chỉ được dùng như thuộc tính bổ sung (\textit{secondary attribute}),
    \item \textbf{ProcessGuid} được sử dụng làm khóa chính (\textit{primary key}) để xây dựng đồ thị tiến trình.
\end{itemize}

\section{Chuẩn hóa dữ liệu và biểu diễn tiến trình}
\label{section:2.2}

\subsection{Chuẩn hóa đường dẫn tiến trình}
\label{subsection:2.2.1}
Trong log \textbf{Windows}, cùng một tiến trình có thể xuất hiện dưới nhiều dạng biểu diễn khác nhau:
\begin{itemize}
    \item \texttt{C:\textbackslash Windows\textbackslash System32\textbackslash cmd.exe}
    \item \texttt{c:\textbackslash windows\textbackslash system32\textbackslash CMD.EXE}
    \item \texttt{\%SystemRoot\%\textbackslash System32\textbackslash cmd.exe}
\end{itemize}

Nếu không chuẩn hóa (\textit{normalization}), hệ thống sẽ coi đây là ba tiến trình khác nhau, làm sai lệch thống kê tần suất và quan hệ tiến trình. Do đó, hệ thống áp dụng các bước chuẩn hóa:

\begin{itemize}
    \item Chuyển toàn bộ ký tự sang \textbf{lowercase},
    \item Chuyển dấu \texttt{\textbackslash} thành \texttt{/},
    \item Mở rộng biến môi trường (\textit{environment variables expansion}) như \texttt{\%SystemRoot\%},
    \item Loại bỏ các biểu diễn dư thừa trong đường dẫn (\textit{path}).
\end{itemize}

Sau khi chuẩn hóa, chuỗi: \textbf{c:/windows/system32/cmd.exe} được sử dụng như biểu diễn duy nhất (\textit{unique representation}).

\subsection{Biểu diễn đồ thị tiến trình}
\label{subsection:2.2.2}
Từ các sự kiện \textbf{Sysmon Event ID 1}, hệ thống xây dựng đồ thị tiến trình có hướng:
\begin{equation}
    G = (V, E)
\end{equation}
trong đó:
\begin{itemize}
    \item Mỗi \textit{node} $v \in V$ là một tiến trình,
    \item Mỗi cạnh $e = (u, v) \in E$ biểu diễn quan hệ sinh tiến trình (\textit{process creation relationship}).
\end{itemize}

Mỗi \textit{node} được gắn các thuộc tính (\textit{attributes}):
\begin{itemize}
    \item \textit{Command line},
    \item \textit{Timestamp},
    \item \textit{Parent process},
    \item \textit{Network activity},
    \item \textbf{Sigma score},
    \item Và các đặc trưng thống kê (\textit{statistical features}) khác.
\end{itemize}

Khác với cây tiến trình (\textit{process tree}) truyền thống, đồ thị này giữ lại toàn bộ quan hệ thực thi, phục vụ cho các bước phân tích chuỗi (\textit{sequence analysis}) và lan truyền rủi ro (\textit{risk propagation}).

\section{Phân loại bất thường theo mô hình Chandola}
\label{subsection:2.3}
Theo Chandola et al. (2009), bất thường được chia thành ba loại chính:

\begin{itemize}
    \item Point Anomalies: một thực thể đơn lẻ có hành vi bất thường, ví dụ command line chứa chuỗi mã hóa base64.
    \item Contextual Anomalies: hành vi bình thường nhưng xuất hiện trong ngữ cảnh không phù hợp, ví dụ word.exe sinh cmd.exe.
    \item Collective Anomalies: một nhóm hành vi bình thường nhưng khi kết hợp lại tạo thành bất thường, ví dụ nhiều tiến trình tạo kết nối mạng trong thời gian ngắn.
\end{itemize}

Phân loại này cung cấp cơ sở để xây dựng các mô hình phát hiện đa chiều, thay vì chỉ dựa vào một đặc trưng đơn lẻ.


\section{Cơ sở toán học của mô hình định lượng rủi ro}
\label{section:2.6}
Trong các hệ thống điều tra số hiện đại, đặc biệt trên nền tảng \textbf{Windows}, bài toán quan trọng không chỉ dừng lại ở việc phát hiện một tiến trình có độc hại hay không, mà còn cần lượng hóa được mức độ nguy hiểm của tiến trình đó dựa trên các tín hiệu quan sát được từ \textit{log} hệ thống. Điều này đặc biệt cần thiết trong môi trường \textbf{SOC} thực tế, nơi số lượng sự kiện có thể đạt tới hàng triệu bản ghi mỗi ngày và việc điều tra thủ công toàn bộ dữ liệu là không khả thi.

Về bản chất, bài toán có thể được mô hình hóa như một hàm ánh xạ:
\begin{equation}
    Score = f(X)
\end{equation}
trong đó $X$ là tập các đặc trưng mô tả tiến trình, bao gồm hành vi nội tại, quan hệ tiến trình cha–con, ngữ cảnh thực thi, tín hiệu từ luật phát hiện và các đặc trưng thống kê khác.

Trong thực tế, các đặc trưng này phản ánh nhiều khía cạnh khác nhau của hành vi bất thường và mang tính độc lập tương đối. Do đó, nếu chỉ sử dụng các phép cực trị như $Score = \max(x_i)$ thì hệ thống chỉ giữ lại tín hiệu mạnh nhất và bỏ qua các tín hiệu hỗ trợ. Điều này không phù hợp với đặc thù của các cuộc tấn công \textit{APT}, nơi hành vi độc hại thường được biểu hiện thông qua sự kết hợp của nhiều tín hiệu mức trung bình thay vì một dấu hiệu duy nhất.

Vì vậy, hệ thống lựa chọn mô hình tổ hợp tuyến tính có trọng số:
\begin{equation}
    Score \approx \sum_{i} w_i x_i
\end{equation}
trong đó $x_i$ là các thành phần bất thường và $w_i$ là trọng số phản ánh tầm quan trọng tương đối của từng thành phần. Dựa trên phân loại của \textbf{Varun Chandola} về phát hiện bất thường, hệ thống chia các tín hiệu thành năm nhóm chính:
\begin{itemize}
    \item \textbf{Bất thường điểm} (\textit{Point Anomaly} – $P$),
    \item \textbf{Bất thường ngữ cảnh} (\textit{Contextual Anomaly} – $C$),
    \item \textbf{Tín hiệu từ luật phát hiện} (\textit{Sigma} – $S$),
    \item \textbf{Bất thường chuỗi} (\textit{Sequence Anomaly} – $Seq$),
    \item \textbf{Bất thường tập thể} (\textit{Collective Anomaly} – $Col$).
\end{itemize}

Từ đó, mô hình tổng quát được xây dựng:
\begin{equation}
    Final\_Score = (w_1P + w_2C + w_3S + w_4Seq) \cdot CI + w_5Col
\end{equation}
trong đó \textbf{CI} là hệ số đồng thuận (\textit{Corroboration Index}), dùng để phản ánh mức độ cộng hưởng giữa các tín hiệu độc lập. Khi nhiều đặc trưng độc lập cùng chỉ ra
bất thường, giá trị này sẽ tăng, giúp nhấn mạnh các trường hợp có độ tin cậy cao.
Ngược lại, nếu chỉ có một tín hiệu đơn lẻ, hệ số này sẽ giảm để hạn chế ảnh hưởng
của nhiễu.


\subsection{Chuẩn hóa tần suất xuất hiện – NIF (Normalized Inverse Frequency)}
\label{subsection:2.4.1}
Trong dữ liệu log \textbf{Windows}, tần suất xuất hiện của tiến trình là một tín hiệu quan trọng phản ánh mức độ phổ biến của hành vi. Tuy nhiên, việc sử dụng trực tiếp tần suất tuyệt đối là không phù hợp do phân phối tiến trình trong hệ thống mang tính lệch mạnh (\textit{heavy-tailed distribution}). Một số tiến trình hệ thống như \textbf{svchost.exe}, \textbf{explorer.exe} hoặc \textbf{chrome.exe} có thể xuất hiện hàng nghìn lần, trong khi các công cụ quản trị hoặc công cụ \textit{recon} như \textbf{whoami.exe}, \textbf{quser.exe} hoặc \textbf{nltest.exe} có thể chỉ xuất hiện một vài lần.

Nếu sử dụng trực tiếp số lần xuất hiện $f_i$, thì các tiến trình hiếm gần như luôn bị đánh giá là bất thường, mặc dù nhiều hành vi hiếm vẫn hoàn toàn hợp lệ trong thực tế. Do đó, cần một đại lượng phản ánh mức độ ``hiếm tương đối'' thay vì tần suất tuyệt đối. Gọi $f_i$ là số lần xuất hiện của tiến trình $image_i$, và $f_{max} = \max_i(f_i)$ là tần suất lớn nhất trong toàn bộ tập dữ liệu.

Hệ thống định nghĩa chuẩn hóa tần suất nghịch đảo (\textbf{NIF} - \textit{Normality Isolation Factor}):
\begin{equation}
    NIF_i = 1 - \frac{\log(f_i)}{\log(f_{max})}
\end{equation}

Việc sử dụng hàm $\log$ có ý nghĩa quan trọng về mặt thống kê. Trong phân phối lệch mạnh, khoảng cách giữa các tiến trình phổ biến và không phổ biến có thể lên tới nhiều bậc độ lớn. Nếu sử dụng tỷ lệ tuyến tính $1 - \frac{f_i}{f_{max}}$, thì phần lớn tiến trình hiếm sẽ bị dồn về giá trị gần 1 và mất khả năng phân biệt. Hàm $\log$ giúp nén khoảng cách này, đồng thời vẫn giữ được tính thứ tự giữa các tiến trình.

Theo định nghĩa trên:
\begin{itemize}
    \item Nếu $f_i = f_{max}$, thì $NIF_i = 0$.
    \item Nếu $f_i = 1$, thì $NIF_i = 1$.
\end{itemize}
Do đó, $NIF_i \in [0, 1]$ và có thể được sử dụng trực tiếp như một đặc trưng chuẩn hóa. 

Ngoài tần suất của tên tiến trình, hệ thống còn áp dụng \textbf{NIF} cho các quan hệ tiến trình cha–con $(parent, child)$ nhằm đánh giá mức độ hiếm của hành vi sinh tiến trình. Điều này đặc biệt hữu ích trong các kịch bản như $\textbf{winword.exe} \rightarrow \textbf{cmd.exe}$ hoặc $\textbf{powershell.exe} \rightarrow \textbf{certutil.exe}$, vốn hiếm gặp trong hoạt động thông thường nhưng xuất hiện phổ biến trong nhiều chiến dịch tấn công.

\subsection{Phát hiện bất thường điểm (Point Anomaly) – Điểm P}
\label{subsection:2.4.2}
Bất thường điểm (\textit{Point Anomaly}) là dạng bất thường cổ điển nhất trong phát hiện dị thường (\textit{anomaly detection}), tương ứng với các điểm dữ liệu nằm cách xa phần lớn quần thể. Trong bối cảnh log \textbf{Windows}, điều này tương ứng với các tiến trình có hành vi khác biệt rõ rệt về mặt thống kê so với phần còn lại của hệ thống.

Điểm bất thường điểm được ký hiệu là $P$ và được xây dựng từ hai thành phần chính:
\begin{itemize}
    \item Bất thường toàn cục (\textit{global anomaly}) dựa trên \textbf{MAD},
    \item Bất thường nội bộ nhóm (\textit{local anomaly}) dựa trên \textbf{IPAS}.
\end{itemize}


\subsubsection{Median Absolute Deviation (MAD)}
\label{subsubsection:2.4.2.1}
Các đặc trưng hành vi như độ dài \textit{command line}, \textit{entropy} chuỗi lệnh, số lượng tiến trình con hoặc số kết nối mạng thường có phân phối lệch và chứa nhiều \textit{outlier}. Trong trường hợp này, việc sử dụng trung bình (\textit{mean}) và độ lệch chuẩn (\textit{standard deviation}) sẽ thiếu ổn định do rất nhạy với các giá trị cực đoan.

Để giải quyết vấn đề này, hệ thống sử dụng \textbf{Median Absolute Deviation} (\textbf{MAD}), một đại lượng thống kê \textit{robust} có khả năng chống nhiễu cao. Cho tập dữ liệu $X = \{x_1, x_2, ..., x_n\}$, gọi $m = \text{median}(X)$ là trung vị của tập dữ liệu. Khi đó, \textbf{MAD} được xác định:
\begin{equation}
    \text{MAD} = \text{median}(|x_i - m|)
\end{equation}

Từ đó, \textbf{robust z-score} được xác định:
\begin{equation}
    z_{robust} = \frac{0.6745(x_i - m)}{\text{MAD}}
\end{equation}

Hệ số $0.6745$ xuất phát từ quan hệ giữa \textbf{MAD} và độ lệch chuẩn trong phân phối chuẩn: $\text{MAD} \approx 0.6745\sigma$. Nhờ đó, \textbf{robust z-score} vẫn có thể được diễn giải tương tự \textbf{z-score} truyền thống nhưng bền vững hơn đối với \textit{outlier}.

Điểm bất thường dựa trên \textbf{MAD} được xây dựng dưới dạng:
\begin{equation}
    P_{MAD} = g(z_{robust})
\end{equation}
trong đó $g(\cdot)$ là hàm \textit{scale} đưa giá trị về khoảng $[0, 10]$. Cách tiếp cận này giúp phản ánh mức độ bất thường theo khoảng cách thống kê thay vì sử dụng phân loại nhị phân ``bình thường/bất thường''.

\subsubsection{IPAS – Intra-binary Population Anomaly Score}
\label{subsubsection:2.4.2.2}
Trong nhiều trường hợp, một tiến trình không bất thường ở mức toàn cục (\textit{global}) nhưng lại khác biệt rõ rệt so với các tiến trình cùng loại. Ví dụ, trong hàng trăm tiến trình \textbf{powershell.exe}, chỉ một tiến trình duy nhất thực hiện \textit{obfuscation} hoặc \textit{injection}. Nếu chỉ sử dụng \textbf{MAD} toàn cục, các trường hợp này có thể bị pha loãng bởi số lượng lớn tiến trình hợp lệ cùng tên.

Để giải quyết vấn đề này, hệ thống sử dụng \textbf{Intra-binary Population Anomaly Score} (\textbf{IPAS}). Đối với mỗi nhóm tiến trình có cùng tên tệp, hệ thống xây dựng vector đặc trưng $V_{10}$ gồm 10 chiều hành vi. Với mỗi tiến trình $j$, khoảng cách tới trung tâm nhóm được tính:

\begin{equation}
    dist_j = \frac{1}{10} \sum_{k=1}^{10} \frac{|v_{j,k} - \mu_k|}{\sigma_k + \epsilon}
\end{equation}

trong đó:
\begin{itemize}
    \item $\mu_k$ là giá trị trung bình của chiều thứ $k$,
    \item $\sigma_k$ là độ lệch chuẩn (\textit{standard deviation}),
    \item $\epsilon$ là hằng số nhỏ để tránh lỗi chia cho 0.
\end{itemize}

Biểu thức này tương đương với một dạng đơn giản hóa của khoảng cách \textbf{Mahalanobis} trong trường hợp giả sử các chiều độc lập. Giá trị khoảng cách càng lớn thì tiến trình càng khác biệt so với các tiến trình cùng loại. Sau đó, điểm \textbf{IPAS} được xây dựng:

\begin{equation}
    IPAS = h(dist_j)
\end{equation}

trong đó $h(\cdot)$ là hàm \textit{scale} đưa giá trị về miền $[0, 10]$.

\subsubsection{Tổng hợp điểm P}
\label{subsubsection:2.4.2.3}
Điểm bất thường điểm cuối cùng được xây dựng bằng cách kết hợp bất thường toàn cục và bất thường nội bộ nhóm:
\begin{equation}
    P = \alpha P_{MAD} + \beta IPAS
\end{equation}
trong đó $\alpha + \beta = 1$. Cách kết hợp này giúp hệ thống vừa phát hiện được các hành vi cực đoan trên toàn bộ tập dữ liệu, vừa giữ được khả năng phát hiện các tiến trình bị ``vũ khí hóa'' (\textit{weaponized}) lẩn khuất trong cùng một nhóm tiến trình phổ biến.

Ngoài ra, hệ thống còn xem xét thông tin từ quá trình phân cụm \textbf{HDBSCAN}. Các tiến trình thuộc \textit{cluster} nhiễu ($cluster = -1$) được xem là các điểm không thuộc bất kỳ cấu trúc mật độ ổn định nào và thường mang tính dị thường cao. Do đó, các tiến trình này sẽ được tăng trọng số bất thường trong quá trình tính điểm.


\subsection{Phát hiện bất thường điểm (Point Anomaly) – Điểm P}
\label{subsection:2.4.3}




%%%%%%%%%%%%%%%%%%%%%%%%%%%%%%%%%%%

\end{document}